\documentclass[11pt,a4paper]{article}

\usepackage[margin=2.5cm]{geometry}
\usepackage{amsmath,amssymb,amsthm}
\usepackage{booktabs}
\usepackage{graphicx}
\usepackage{tikz}
\usetikzlibrary{arrows.meta,positioning}
\usepackage[T1]{fontenc}
\usepackage[utf8]{inputenc}
\usepackage[expansion=false]{microtype}
\usepackage{array}
\usepackage{algorithm}
\usepackage{algpseudocode}
\usepackage{hyperref}

\theoremstyle{definition}
\newtheorem{definition}{Definition}
\theoremstyle{plain}
\newtheorem{theorem}{Theorem}
\newtheorem{lemma}[theorem]{Lemma}
\newtheorem{proposition}[theorem]{Proposition}
\newtheorem{corollary}[theorem]{Corollary}
\theoremstyle{remark}
\newtheorem{fact}{Fact}
\newtheorem{remark}{Remark}
\newtheorem{example}{Example}

\newcommand{\cov}{\operatorname{cov}}
\newcommand{\orig}{\operatorname{orig}}
\newcommand{\name}{\operatorname{name}}
\newcommand{\isize}{\operatorname{isize}}
\newcommand{\Feas}{\mathcal{F}}
\newcommand{\Links}{\mathcal{K}}
\newcommand{\Names}{\mathcal{N}}
\newcommand{\scap}{\operatorname{scap}}

\algrenewcommand\algorithmicrequire{\textbf{Input:}}
\algrenewcommand\algorithmicensure{\textbf{Output:}}

\title{\texttt{tree-dp3}: exact cluster-local re-optimization\\
for differential compression of the gapped suffix array}
\author{}
\date{}

\begin{document}
\maketitle

\begin{abstract}
\noindent
The differential compression of a gapped suffix array chooses a set of
\emph{links}, each replacing a run of stored positions by a pointer plus
a shift. We show that the resulting optimization problem is a
maximum-weight set packing subject to exactly two constraints --- at most
one link per name, and a chosen link's source ranks must all be kept ---
and that, crucially, \emph{both are pairwise} and there is no acyclicity
requirement, because the compressed array stores literal positions and
the decoder never recurses. Pairwise feasibility licenses an exact
\emph{local} optimization: freeze the links of all but a bounded set of
names and re-solve that set optimally. We prove the composition theorem
that makes this sound --- a link's covered ranks always lie in its own
interval, so a frozen link can never drop a rank belonging to a cluster
member, and the entire interface to the frozen part collapses to two
conditions --- and identify the correct notion of dependency between
names (one's link sources from the other's interval), showing why the two
tempting alternatives (competition for covered ranks, sharing a source)
are not conflicts. The resulting algorithm, \texttt{tree-dp3}, runs the
existing preference-forest DP and unpin/retarget loop and then performs
branch-and-bound over bounded dependency clusters. It repairs both known
counterexamples to the forest DP, attains the ILP optimum on $9$ of $10$
fixed test instances, and on a $139$-configuration suite at
$a_{\max}=16$ reduces total $|C|$ from $252329$ to $213078$
($15.5\,\%$), improving $122$ configurations and worsening none. At the
default $a_{\max}=8$ the same exact solve buys only $3$ positions
suite-wide: the binding constraint is the candidate universe, not the
neighbourhood size.
\end{abstract}

\tableofcontents

\section{The compression problem}
\label{sec:problem}

\subsection{Setting and notation}

We reuse the setting of the walkthrough (\texttt{dep\_order\_walkthrough})
and of the pseudoforest note (\texttt{exact\_pseudoforest\_dp}) without
change. Fix a binary shape $S$ of span $s$ over a text $T[0..n)$. The
DisLex transform lays the gapped $k$-mer names out grouped by residue
class $p \bmod s$; the suffix array of the resulting lextext is the
\textbf{gapped suffix array} (GSA), with $m=n+1$ ranks. We write
$\orig(r)$ for the original text position of rank $r$, and $\name(r)$ for
the gapped $k$-mer at $\orig(r)$, which is the first symbol of the
rank-$r$ gapped suffix.

Because the GSA is sorted by gapped suffix, the ranks carrying a fixed
name $c$ form a contiguous \textbf{interval}
\[
  I_c = [\ell_c, r_c) \subseteq \{0,\dots,m-1\},
  \qquad
  \isize(c) = |I_c| = r_c - \ell_c .
\]
The intervals $\{I_c\}_c$ \emph{partition} $\{0,\dots,m-1\}$: every rank
has exactly one first symbol. We write $\Names$ for the set of names $c$
with $\isize(c) > 1$; singletons are never compressed. Throughout,
$a_{\max}$ is the maximum shift considered (\texttt{--max-add}, default
$8$) and $\kappa = 3$ is the minimum coverage
(\texttt{kMinCoverage}).

\subsection{What is stored, and what a link is}

Compression drops a subset of ranks. The surviving positions, in rank
order, form the \textbf{compressed array} $C$. For each name $c$ the hash
table stores at most two entries,
\[
  H_{\mathrm{offset}} = (\mathit{pos}, a, \mathit{num}),
  \qquad
  H_{\mathrm{rest}} = (\mathit{pos}', \mathit{num}'),
\]
decoding to
\begin{equation}
\label{eq:decode}
  \text{positions}(c)
  \;=\;
  \bigl\{\, C[\mathit{pos}+t] + a\cdot s \;:\; 0 \le t < \mathit{num} \,\bigr\}
  \;\cup\;
  \bigl\{\, C[\mathit{pos}'+t] \;:\; 0 \le t < \mathit{num}' \,\bigr\}.
\end{equation}
Everything the compressor decides is which $H_{\mathrm{offset}}$ entries
to install. We call one such decision a \emph{link}.

\begin{definition}[Link]
\label{def:link}
A \textbf{link} is a quadruple $k = (c_k,\, a_k,\, S_k,\, T_k)$ where
$c_k$ is a name, $1 \le a_k \le a_{\max}$ is the \emph{add},
$S_k = [\sigma_k, \tau_k)$ is a contiguous range of ranks (the
\textbf{source run}), and
$T_k = (t_0,\dots,t_{|S_k|-1})$ is the sequence of \textbf{covered ranks},
subject to
\begin{enumerate}
\item[(V1)] for every $0 \le j < |S_k|$:
  $\orig(t_j) = \orig(\sigma_k + j) + a_k\cdot s$ and $\name(t_j) = c_k$
  --- so $T_k \subseteq I_{c_k}$, and the map $\sigma_k+j \mapsto t_j$ is
  a bijection $S_k \to T_k$ that is order-preserving;
\item[(V2)] $S_k \cap T_k = \emptyset$;
\item[(V3)] $\cov(k) := |T_k| = |S_k| \ge \kappa$.
\end{enumerate}
We write $\Links$ for the set of all links and
$\Links_c = \{k \in \Links : c_k = c\}$.
\end{definition}

Condition (V1) is the note's order-preservation argument: prepending
$a$ symbols to a gapped suffix preserves relative order and shifts the
original position by exactly $+a\cdot s$. It is checked in the code by
\texttt{pred\_lexpos\_} / \texttt{succ\_lexpos\_}. Condition (V2) is an
intrinsic validity condition, not an interaction between two links: a
link may not use as a source a rank it itself drops. It is what
\texttt{emit\_intra\_runs\_} enforces by a two-pointer scan, and what the
\texttt{clash} test of \texttt{enumerate\_all\_links\_} enforces in the
source-centric enumeration. Note that (V2) does \emph{not} require
$S_k \cap I_{c_k} = \emptyset$: intra-interval links, whose source run
overlaps the target's own interval but avoids the ranks it covers, are
perfectly legal, and \texttt{GCSA\_INTRA\_LINKS=0} exists only to restore
the historical, stricter rule.

\subsection{The objective and the two constraints}

Let $X \subseteq \Links$ be a set of accepted links. A rank is dropped
from $C$ exactly when it is covered by some $k \in X$. Hence

\begin{equation}
\label{eq:objective}
  |C|(X) \;=\; m \;-\; \cov(X),
  \qquad
  \cov(X) := \sum_{k \in X} \cov(k),
\end{equation}
\emph{provided} the covered sets are disjoint --- which
Proposition~\ref{prop:target-free} below shows is automatic. Minimizing
$|C|$ is therefore maximizing total coverage.

\begin{definition}[Feasibility]
\label{def:feasible}
$X \subseteq \Links$ is \textbf{feasible} iff
\begin{itemize}
\item[(F1)] \textbf{one link per name}: $c_k \ne c_{k'}$ for all
  distinct $k,k' \in X$;
\item[(F2)] \textbf{source integrity}: $S_k \cap T_{k'} = \emptyset$ for
  all $k, k' \in X$ with $k \ne k'$.
\end{itemize}
We write $\Feas$ for the family of feasible sets.
\end{definition}

(F1) is a storage restriction: a \texttt{HashEntry} holds a single
$H_{\mathrm{offset}}$. (F2) says a source must be physically present in
$C$: if some other accepted link dropped a rank of $S_k$, then
$C[\mathit{pos}+t]$ in \eqref{eq:decode} would silently read a different
rank's position. In the implementation, (F2) is what
\texttt{run\_kept\_} verifies before an accept, and what the
\texttt{pin\_count\_} array tracks incrementally.

\begin{proposition}[Target-disjointness is free]
\label{prop:target-free}
If $X$ satisfies (F1) then the covered sets $\{T_k\}_{k \in X}$ are
pairwise disjoint, and \eqref{eq:objective} holds exactly.
\end{proposition}

\begin{proof}
Let $k \ne k' \in X$. By (F1), $c_k \ne c_{k'}$. By (V1),
$T_k \subseteq I_{c_k}$ and $T_{k'} \subseteq I_{c_{k'}}$, and distinct
names have disjoint intervals. Hence $T_k \cap T_{k'} = \emptyset$, so
$|\bigcup_{k \in X} T_k| = \sum_{k\in X}\cov(k)$.
\end{proof}

This small observation is the seed of everything that follows: because a
link's covered ranks are confined to its own interval, and (F1) makes the
per-name decision unique, \emph{the name owns its ranks}. We will use
this twice more, in Lemma~\ref{lem:dep} and Theorem~\ref{thm:compose}.

\subsection{There is no acyclicity constraint}
\label{sec:noacyclic}

It is tempting to expect a third constraint of the form ``the
points-to relation among names must be acyclic'', by analogy with
delta-encoded or recursively-defined structures. There is none.

\begin{proposition}[Non-recursive decoding]
\label{prop:noacyclic}
For any feasible $X$, the decoder \eqref{eq:decode} is well defined and
terminates in $O(\mathit{num}+\mathit{num}')$ time for every name,
independently of the structure of $X$.
\end{proposition}

\begin{proof}
$C$ stores \emph{literal} text positions of kept ranks in rank order.
Given $k \in X$, (F2) guarantees every rank of $S_k$ is kept, so $S_k$
occupies $|S_k|$ consecutive entries of $C$ starting at some index
$\mathit{pos}$; this index is exactly what is written into
$H_{\mathrm{offset}}$. Evaluating \eqref{eq:decode} therefore reads
$\mathit{num}$ consecutive entries of $C$ and adds $a\cdot s$ to each. No
term on the right-hand side of \eqref{eq:decode} refers to
$\text{positions}(c')$ for any name $c'$; the definition is not
recursive, so there is nothing to be well-founded.
\end{proof}

A hypothetical scheme in which $H_{\mathrm{offset}}$ could point at
\emph{dropped} (virtual) ranks would decode a chain
$c_1 \to c_2 \to \cdots$, and would need that chain to terminate --- an
acyclicity condition, and a global, non-pairwise one. Constraint (F2)
forbids such chains outright, and in exchange the feasible region stays
simple:

\begin{lemma}[Pairwise feasibility]
\label{lem:pairwise}
$X \in \Feas$ if and only if every $k \in X$ satisfies (V1)--(V3) and
every pair $\{k,k'\} \subseteq X$ with $k \ne k'$ satisfies (F1) and
(F2). Consequently $\Feas$ is downward closed: $X \in \Feas$ and
$Y \subseteq X$ imply $Y \in \Feas$.
\end{lemma}

\begin{proof}
Both (F1) and (F2) are universally quantified over ordered pairs of
distinct elements of $X$ and involve no other member of $X$; so the
conjunction over $X$ is exactly the conjunction over its $2$-element
subsets. Downward closure is immediate, since a conjunction over the
pairs of $Y$ is a sub-conjunction of the one over the pairs of $X$.
\end{proof}

\begin{corollary}[Set-packing formulation]
\label{cor:packing}
Differential compression is the maximum-weight independent set problem
\[
  \min_{X \in \Feas} |C|(X)
  \;=\;
  m - \max_{X \text{ independent in } G} \ \sum_{k \in X} \cov(k),
\]
on the conflict graph $G = (\Links, E)$ with
$\{k,k'\} \in E$ iff $c_k = c_{k'}$, or
$S_k \cap T_{k'} \ne \emptyset$, or $S_{k'} \cap T_k \ne \emptyset$.
\end{corollary}

Corollary~\ref{cor:packing} is the same formulation the pseudoforest note
starts from; what is new here is the emphasis on Lemma~\ref{lem:pairwise}.
Pairwise-ness is precisely the licence for \emph{cluster-local exact
optimization}: if feasibility were a global property of $X$ (an
acyclicity or an ordering condition), freezing part of a solution and
re-solving the rest could produce a globally infeasible result even when
every local check passed. Because it is pairwise, the interaction between
a frozen part and a re-solved part is fully described by a finite set of
two-element checks, and those checks can be precomputed once from the
frozen part --- which is exactly what \S\ref{sec:compose} does.

\section{Background: the forest DP and why it is not optimal}
\label{sec:treedp}

\subsection{Recap of \texttt{tree-dp}}

\texttt{tree-dp} (documented in \texttt{dep\_order\_walkthrough}, \S7)
proceeds as follows.

\begin{enumerate}
\item \textbf{Preferences.} For each name $c$ with $\isize(c) > 2$,
  enumerate $\hat\Links_c$ --- the heuristics' view of the link universe
  (\texttt{enumerate\_candidates\_}: maximal source runs with internal
  $\mathrm{lcp} \ge a+1$) --- and record the \textbf{preferred} link
  $\pi(c)$, of maximum coverage with ties broken towards larger $a$.
  Preferences ignore availability; they are a static, structural notion.
\item \textbf{Preference graph.} Draw a directed edge $c \to w$ whenever
  $\pi(c)$'s source lies inside $I_w$. Since $S_{\pi(c)}$ has
  $\mathrm{lcp} \ge a+1 \ge 2$ throughout, it lies inside a single
  interval, so this graph has out-degree $\le 1$: as an undirected graph
  it is a pseudoforest.
\item \textbf{Cycle-breaking.} The implementation processes preferences
  in ascending name order and installs $\mathrm{parent}(c) := w$ only if
  $c$ is not already an ancestor of $w$ (\texttt{has\_ancestor}). Edges
  that would close a cycle are \emph{discarded}. The result is a forest.
\item \textbf{KEEP/COMPRESS DP.} With
  $\mathit{src\_ok}$ recording whether the node's preferred source is
  still available,
  \[
    \mathrm{DP}(v,\mathsf{true}) = \min\Bigl(
      \isize(v) + \!\!\sum_{w \in \mathrm{ch}(v)}\!\! \mathrm{DP}(w,\mathsf{true}),\;\;
      \isize(v) - \cov(\pi(v)) + \!\!\sum_{w \in \mathrm{ch}(v)}\!\! \mathrm{DP}(w,\mathsf{false})
    \Bigr),
  \]
  \[
    \mathrm{DP}(v,\mathsf{false}) = \isize(v) + \sum_{w \in \mathrm{ch}(v)} \mathrm{DP}(w,\mathsf{true}),
  \]
  with KEEP on ties. A root may COMPRESS only when its preferred source
  lies outside its own tree.
\item \textbf{Materialize, leftover, Phase II.} The DP's choices are
  accepted deepest-first; names that never entered the forest are handled
  by a size-ordered availability greedy; then the shared Phase II
  unpin/retarget fixed point runs.
\end{enumerate}

The DP is exact \emph{for the model it is given}. Two restrictions of
that model cost optimality: each name is allowed only its single
preferred link, and each name is allowed only its single forest parent,
with any surplus preference edge silently deleted in step~3. The two
counterexamples below isolate the two failures.

\subsection{Counterexample I: cycle-breaking fixes the wrong orientation}
\label{sec:ce1}

Take $T = (\mathtt{AC})^6 = \mathtt{ACACACACACAC}$, $n = 12$,
$S = \mathtt{\#.\#}$, $s = 3$, $m = 13$. Position $p$ carries the name
$\mathtt{AA}$ for even $p \le 8$ and $\mathtt{CC}$ for odd $p \le 9$; the
three tail positions give the singletons $\mathtt{A\$}, \mathtt{C\$},
\mathtt{\$\$}$. The GSA is

\begin{center}
\small
\begin{tabular}{@{}lrrrrrrrrrrrrr@{}}
\toprule
rank $r$ & 0 & 1 & 2 & 3 & 4 & 5 & 6 & 7 & 8 & 9 & 10 & 11 & 12 \\
$\orig(r)$ & 12 & 10 & 8 & 6 & 4 & 2 & 0 & 11 & 9 & 7 & 5 & 3 & 1 \\
$\name(r)$ & \$\$ & A\$ & AA & AA & AA & AA & AA & C\$ & CC & CC & CC & CC & CC \\
\bottomrule
\end{tabular}
\end{center}

so $I_{\mathtt{AA}} = [2,7)$ and $I_{\mathtt{CC}} = [8,13)$, both of size
$5$, and $m = 13 = 5 + 5 + 3$.

\paragraph{The whole link universe.}
A link shifts original positions by $+a\cdot s = 3a$. Since AA sits on
even and CC on odd positions, $3a$ must be odd, i.e.\ $a$ odd, and
$a \ge 3$ overshoots the text. So $a = 1$ throughout, and the universe
(as reported by \texttt{ilp\_baseline}: \texttt{candidates=4}) is

\begin{center}
\small
\begin{tabular}{@{}llllrl@{}}
\toprule
link & target & $a$ & source $S_k$ & $\cov$ & covered $T_k$ \\
\midrule
$k_1$   & AA & 1 & $[10,13)$, i.e.\ $\orig = 5,3,1$ & 3 & $(2,3,4)$, i.e.\ $\orig = 8,6,4$ \\
$k_2$   & CC & 1 & $[3,7)$,\ \ i.e.\ $\orig = 6,4,2,0$ & 4 & $(8,9,10,11)$, i.e.\ $\orig = 9,7,5,3$ \\
$k_2'$  & CC & 1 & $[3,6)$ & 3 & $(8,9,10)$ \\
$k_2''$ & CC & 1 & $[4,7)$ & 3 & $(9,10,11)$ \\
\bottomrule
\end{tabular}
\end{center}

($k_2', k_2''$ are the two sub-runs of $k_2$ of length $\ge\kappa$; only
$k_1$ and $k_2$ are maximal, hence only those two are visible to
\texttt{enumerate\_candidates\_}.)

\paragraph{Every feasible set has coverage at most $4$.}
$S_{k_1} = \{10,11,12\}$ meets $T_{k_2} = \{8,9,10,11\}$ in $\{10,11\}$,
meets $T_{k_2'}$ in $\{10\}$ and $T_{k_2''}$ in $\{10,11\}$; so $k_1$
conflicts with every CC link by (F2). Symmetrically
$S_{k_2} = \{3,4,5,6\}$ meets $T_{k_1} = \{2,3,4\}$ in $\{3,4\}$. At most
one of the two names can compress, and the best single choice is $k_2$:
\[
  |C|^\star = 13 - 4 = 9,
\]
which \texttt{ilp\_baseline} confirms (\texttt{optimal |C| = 9},
$1$ selected candidate).

\paragraph{What \texttt{tree-dp} does.}
The preferred links are
$\pi(\mathtt{AA}) = k_1$ ($\cov 3$) and $\pi(\mathtt{CC}) = k_2$
($\cov 4$), giving a $2$-cycle
$\mathtt{AA} \leftrightarrow \mathtt{CC}$ in the preference graph. Both
names have $\isize = 5 > 2$ and both preferences clear $\cov \ge \kappa$,
so, unlike on the $\mathtt{GCCTTTAAAG}^3$ instance of the walkthrough,
the cycle is \emph{not} filtered away and the cycle-breaking rule fires.
Names are processed in ascending order, and $\mathtt{AA} < \mathtt{CC}$
as arithmetic codes ($2 < 8$), so:

\begin{center}
\small
\begin{tabular}{@{}llp{9.2cm}@{}}
\toprule
step & edge & outcome \\
\midrule
1 & $\mathtt{AA} \to \mathtt{CC}$ &
  installed: $\mathrm{parent}(\mathtt{AA}) = \mathtt{CC}$ \\
2 & $\mathtt{CC} \to \mathtt{AA}$ &
  \textbf{discarded}: \texttt{has\_ancestor(AA,\,CC)} is true, so this
  edge would close a cycle \\
\bottomrule
\end{tabular}
\end{center}

The forest is the single edge $\mathtt{CC} \to \mathtt{AA}$ with root
$\mathtt{CC}$. The DP then evaluates
\[
\begin{aligned}
  \mathrm{DP}(\mathtt{AA},\mathsf{true}) &= \min(5,\ 5-3) = 2 \quad(\text{COMPRESS}),\\
  \mathrm{cost}_{\mathrm{keep}}(\mathtt{CC}) &= 5 + 2 = 7,
\end{aligned}
\]
and COMPRESS is \emph{not available} at the root: $\mathtt{CC}$'s
preferred source $\mathtt{AA}$ is a forest descendant, which is exactly
the special case the root rule patches. So $\mathtt{CC}$ is forced to
KEEP and the DP returns $7$, i.e.
\[
  |C| \;=\; 7 + 3\ \text{singletons} \;=\; 10 .
\]
The subsequent leftover greedy and Phase II find nothing (both names are
saturated), so $\texttt{tree-dp} = \texttt{tree-dp2} = 10$ against the
optimum $9$.

\paragraph{Diagnosis.}
The instance contains a genuine, high-coverage mutual preference, and the
two orientations are \emph{not} equally good: $\mathtt{CC}\!\leftarrow\!\mathtt{AA}$
is worth $4$, $\mathtt{AA}\!\leftarrow\!\mathtt{CC}$ only $3$.
Cycle-breaking does not compare them; it keeps whichever edge it meets
first in a deterministic but semantically arbitrary order, and the DP
then optimizes exactly --- inside a model from which the better option
has already been removed. The pseudoforest DP of
\texttt{exact\_pseudoforest\_dp} fixes precisely this failure by folding
the cycle instead of cutting it. Note that plain \texttt{dep-order}
happens to get $9$ here (its Phase II retarget lets $\mathtt{CC}$ take
the coverage-$4$ link), while size-order greedy gets $10$: no method
below is uniformly better than another, which is itself part of the
motivation for an exact local solve.

\subsection{Counterexample II: the optimum is not built from preferred links}
\label{sec:ce2}

Take $T = (\mathtt{ACGT})^5$, $n = 20$, $S = \mathtt{\#.\#}$, $m = 21$.
The four compressible names and their preferred links are

\begin{center}
\small
\begin{tabular}{@{}lrllrl@{}}
\toprule
name $c$ & $\isize$ & $I_c$ & $\pi(c)$ sources from & $\cov$ & detail \\
\midrule
AG & 5 & $[1,6)$   & CT & 4 & $a=1$, $S = [7,11)$,  $T = (1,2,3,4)$ \\
CT & 5 & $[6,11)$  & GA & 4 & $a=1$, $S = [12,16)$, $T = (6,7,8,9)$ \\
GA & 4 & $[12,16)$ & AG & 3 & $a=2$, $S = [3,6)$,   $T = (12,13,14)$ \\
TC & 4 & $[17,21)$ & AG & 4 & $a=1$, $S = [2,6)$,   $T = (17,18,19,20)$ \\
\bottomrule
\end{tabular}
\end{center}

plus three singletons ($\mathtt{\$\$}$, $\mathtt{G\$}$, $\mathtt{T\$}$),
so $m = 5+5+4+4+3 = 21$.

The preference graph has a $3$-cycle
$\mathtt{AG} \to \mathtt{GA} \to \mathtt{CT} \to \mathtt{AG}$
(reading ``$w \to c$'' as ``$c$ prefers a source in $I_w$'') with
$\mathtt{TC}$ hanging off $\mathtt{AG}$. Cycle-breaking in ascending name
order installs $\mathrm{parent}(\mathtt{AG}) = \mathtt{CT}$, then
$\mathrm{parent}(\mathtt{CT}) = \mathtt{GA}$, then \emph{discards}
$\mathtt{GA} \to \mathtt{AG}$ (which would close the cycle), then
installs $\mathrm{parent}(\mathtt{TC}) = \mathtt{AG}$. The forest is the
path
\[
  \mathtt{GA} \;\to\; \mathtt{CT} \;\to\; \mathtt{AG} \;\to\; \mathtt{TC},
\]
rooted at $\mathtt{GA}$, exactly as \texttt{GCSA\_TRACE\_DP=1} reports.
The DP unrolls to
\[
\begin{aligned}
  \mathrm{DP}(\mathtt{TC},\mathsf{true}) &= \min(4,\ 4-4) = 0,
    &\qquad \mathrm{DP}(\mathtt{TC},\mathsf{false}) &= 4, \\
  \mathrm{DP}(\mathtt{AG},\mathsf{true}) &= \min(5+0,\ 1+4) = 5,
    &\qquad \mathrm{DP}(\mathtt{AG},\mathsf{false}) &= 5+0 = 5, \\
  \mathrm{DP}(\mathtt{CT},\mathsf{true}) &= \min(5+5,\ 1+5) = 6,
    &\qquad \mathrm{cost}(\mathtt{GA}) &= 4 + 6 = 10,
\end{aligned}
\]
where $\mathtt{AG}$ takes KEEP on the tie, $\mathtt{CT}$ takes COMPRESS,
and the root $\mathtt{GA}$ has no COMPRESS option at all because its
preferred source $\mathtt{AG}$ is a forest descendant. The DP therefore
accepts two links --- $\mathtt{TC}\!\leftarrow\!\mathtt{AG}$
($\cov 4$) and $\mathtt{CT}\!\leftarrow\!\mathtt{GA}$ ($\cov 4$) --- for
total coverage $8$ and
\[
  |C| \;=\; 21 - 8 \;=\; 13 .
\]

The ILP optimum is $11$, i.e.\ coverage $10$, realised by \emph{three}
links, all sourcing from the (fully kept) interval $I_{\mathtt{AG}}$:

\begin{center}
\small
\begin{tabular}{@{}llrrll@{}}
\toprule
target & source & $a$ & $\cov$ & source ranks & preferred? \\
\midrule
TC & AG & 1 & 4 & $[2,6)$ & yes \\
GA & AG & 2 & 3 & $[3,6)$ & yes --- but this is the edge cycle-breaking deleted \\
CT & AG & 3 & 3 & $[3,6)$ & \textbf{no} ($\pi(\mathtt{CT})$ is
  $\mathtt{CT}\!\leftarrow\!\mathtt{GA}$ at $a=1$, $\cov 4$) \\
\bottomrule
\end{tabular}
\end{center}

Two things go wrong for the forest model simultaneously. First, one of the
three links is on the cycle edge that step~3 discarded. Second, and more
fundamentally, the optimal solution asks $\mathtt{CT}$ to take a
\emph{worse} link than its preferred one ($\cov 3$ at $a=3$ instead of
$\cov 4$ at $a=1$), because doing so moves its source off $I_{\mathtt{GA}}$
and onto $I_{\mathtt{AG}}$, freeing $\mathtt{GA}$ to compress in turn. No
amount of cycle-folding recovers this: the model that allows one link per
name is not rich enough, and a genuinely non-preferred add must be
considered.

Note also that the three optimal links all read from the same ranks
$[3,6) \subseteq S_{\mathtt{TC}} = [2,6)$. That is not an accident, and it
is the phenomenon Lemma~\ref{lem:dep}\,(b) below is about: sources are
shared for free.

\section{The dependency relation}
\label{sec:dep}

We now identify exactly when two names can constrain one another. This is
the relation whose connected balls will become the clusters.

\begin{definition}[Dependency]
\label{def:dep}
Two distinct names $c, d \in \Names$ are \textbf{dependent}, written
$c \sim d$, iff some link of one draws its source from the other's
interval:
\[
  c \sim d
  \quad:\Longleftrightarrow\quad
  \bigl(\exists\, k \in \hat\Links_c :\ S_k \cap I_d \ne \emptyset\bigr)
  \ \ \text{or}\ \
  \bigl(\exists\, k \in \hat\Links_d :\ S_k \cap I_c \ne \emptyset\bigr).
\]
The \textbf{dependency graph} is $D = (\Names, {\sim})$.
\end{definition}

In the implementation this graph is built by walking, for every cached
candidate of every name $c$, the ranks of its source run and adding the
undirected edge $\{c, \name(r)\}$ (deduplicated, and with consecutive
equal names skipped, since a source run typically lies in one interval).

\begin{lemma}[$\sim$ is the only binding relation]
\label{lem:dep}
Let $c \ne d$ be names with $c \not\sim d$. Then for every
$k \in \hat\Links_c$ and every $k' \in \hat\Links_d$, the pair
$\{k,k'\}$ is feasible. Consequently, in the conflict graph of
Corollary~\ref{cor:packing}, there is no edge between $\hat\Links_c$ and
$\hat\Links_d$.
\end{lemma}

\begin{proof}
(F1) cannot bind, since $c_k = c \ne d = c_{k'}$. For (F2), consider
$S_k \cap T_{k'}$. By (V1), $T_{k'} \subseteq I_{d}$, and by
$c \not\sim d$ we have $S_k \cap I_{d} = \emptyset$; hence
$S_k \cap T_{k'} = \emptyset$. The symmetric argument, using
$T_k \subseteq I_c$ and $S_{k'} \cap I_c = \emptyset$, gives
$S_{k'} \cap T_k = \emptyset$.
\end{proof}

Two other relations look like conflicts and are not. Getting either one
wrong would silently enlarge the clusters (making the algorithm slower
without making it better) or, worse, suggest constraints that do not
exist.

\begin{lemma}[Two non-conflicts]
\label{lem:nonconflicts}
Let $k, k'$ be links with $c_k \ne c_{k'}$.
\begin{enumerate}
\item[(a)] \emph{Names never compete for covered ranks:}
  $T_k \cap T_{k'} = \emptyset$ always.
\item[(b)] \emph{Sharing a source is not a conflict:} $S_k \cap S_{k'}$
  may be arbitrary --- in particular $S_k = S_{k'}$ is allowed --- without
  affecting feasibility of $\{k,k'\}$.
\end{enumerate}
\end{lemma}

\begin{proof}
(a) is Proposition~\ref{prop:target-free}: $T_k \subseteq I_{c_k}$,
$T_{k'} \subseteq I_{c_{k'}}$, and distinct names have disjoint
intervals. Note this holds even for intra-interval links, whose
\emph{sources} may straddle interval boundaries; only the covered set is
confined.

(b) Definition~\ref{def:feasible} constrains $S_k$ against $T_{k'}$ and
$S_{k'}$ against $T_k$; the pair $(S_k, S_{k'})$ does not occur in either
(F1) or (F2). Semantically: an accepted link only ever \emph{reads} the
entries of $C$ indexed by its source run, and reading is not exclusive.
This is why \texttt{pin\_count\_} is a counter rather than a lock ---
a rank may be pinned by many links at once, and \texttt{revoke\_} merely
decrements.
\end{proof}

\begin{remark}
Counterexample~II (\S\ref{sec:ce2}) is a witness for
Lemma~\ref{lem:nonconflicts}\,(b) with three links, not two: the optimal
solution has $\mathtt{TC}$, $\mathtt{GA}$ and $\mathtt{CT}$ all sourcing
from $I_{\mathtt{AG}}$, with $S_{\mathtt{GA}} = S_{\mathtt{CT}} = [3,6)$
literally equal and $S_{\mathtt{TC}} = [2,6)$ a superset. Had we modelled
source-sharing as a conflict, that solution would have been declared
infeasible and the algorithm would have reproduced \texttt{tree-dp}'s
$13$.
\end{remark}

\begin{remark}[$\sim$ is necessary but not sufficient]
Definition~\ref{def:dep} quantifies over the whole candidate list of each
name, not over the chosen link, so $c \sim d$ does not imply that the
\emph{selected} links conflict. This is deliberate and conservative: the
cluster must contain every name whose choice could possibly interact with
a member's choice, so $\sim$ over-approximates. Real conflicts are then
tested exactly, per pair of concrete links, during the enumeration
(\S\ref{sec:algo}).
\end{remark}

\section{The composition theorem}
\label{sec:compose}

We can now state the result that makes exact local re-optimization
sound. Fix a feasible $X$ and a set of names $P \subseteq \Names$ (the
\textbf{cluster}). Write
\[
  X_P = \{k \in X : c_k \in P\},
  \qquad
  F = X \setminus X_P \quad (\text{the \textbf{frozen} links}),
\]
and define the two \emph{interface sets} of the frozen part,
\[
  \mathrm{Drop}(F) = \bigcup_{f \in F} T_f,
  \qquad
  \mathrm{Src}(F) = \bigcup_{f \in F} S_f .
\]

\begin{lemma}[Cluster ownership]
\label{lem:ownership}
Let $R_P = \bigcup_{c \in P} I_c$ be the ranks owned by the cluster.
Then $\mathrm{Drop}(F) \cap R_P = \emptyset$. That is, no frozen link can
drop a rank of a cluster member's interval; whether a rank of $R_P$ is
kept or dropped is decided entirely by the links chosen for $P$.
\end{lemma}

\begin{proof}
Let $f \in F$. By definition of $F$, $c_f \notin P$. By (V1),
$T_f \subseteq I_{c_f}$. Since the intervals partition the ranks and
$c_f \ne c$ for every $c \in P$, we get
$T_f \cap I_c = \emptyset$ for every $c \in P$, hence
$T_f \cap R_P = \emptyset$.
\end{proof}

\begin{theorem}[Composition]
\label{thm:compose}
Let $X \in \Feas$, $P \subseteq \Names$, $F = X \setminus X_P$ as above.
Let $Y$ be any set of links such that
\begin{enumerate}
\item[(C1)] every $y \in Y$ satisfies (V1)--(V3) and $c_y \in P$, and
  $Y$ contains at most one link per name;
\item[(C2)] every pair of distinct $y, y' \in Y$ satisfies (F2);
\item[(C3)] $S_y \cap \mathrm{Drop}(F) = \emptyset$ for every $y \in Y$
  \quad (``do not source from a rank a frozen link drops'');
\item[(C4)] $T_y \cap \mathrm{Src}(F) = \emptyset$ for every $y \in Y$
  \quad (``do not drop a rank a frozen link uses as a source'').
\end{enumerate}
Then $F \cup Y \in \Feas$ and
\[
  |C|(F \cup Y) \;=\; |C|(X) \;+\; \cov(X_P) \;-\; \cov(Y).
\]
Conversely, every feasible $Z \in \Feas$ with $Z \setminus Z_P = F$
satisfies (C1)--(C4) with $Y = Z_P$. Hence (C1)--(C4) characterise
exactly the completions of $F$, and
\[
  \min\bigl\{\, |C|(Z) \;:\; Z \in \Feas,\ Z \setminus Z_P = F \,\bigr\}
  \;=\;
  |C|(X) + \cov(X_P) - \max\bigl\{\, \cov(Y) : Y \text{ satisfies (C1)--(C4)} \,\bigr\}.
\]
\end{theorem}

\begin{proof}
\emph{Feasibility.} We check (F1) and (F2) for $F \cup Y$ by
Lemma~\ref{lem:pairwise}, i.e.\ pair by pair. Pairs inside $F$ are
feasible because $F \subseteq X \in \Feas$ and $\Feas$ is downward
closed. Pairs inside $Y$ are (C1) and (C2). For a mixed pair
$f \in F$, $y \in Y$: (F1) holds because $c_f \notin P \ni c_y$; (F2)
requires $S_y \cap T_f = \emptyset$, which is (C3), and
$S_f \cap T_y = \emptyset$, which is (C4).

\emph{Objective.} By Proposition~\ref{prop:target-free} applied to the
feasible sets $X$ and $F \cup Y$, the coverages add, so
$|C|(F\cup Y) = m - \cov(F) - \cov(Y)$ and
$|C|(X) = m - \cov(F) - \cov(X_P)$; subtracting gives the claim.

\emph{Converse.} Let $Z \in \Feas$ with $Z \setminus Z_P = F$ and set
$Y = Z_P$. (C1) is (V1)--(V3) plus (F1) restricted to $Z_P$; (C2) is
(F2) restricted to $Z_P$; and (C3), (C4) are (F2) applied to the mixed
pairs of $Z$, since $\mathrm{Drop}(F)$ and $\mathrm{Src}(F)$ are unions
over $F$ and (F2) is quantified pairwise.
\end{proof}

Theorem~\ref{thm:compose} says the interface between a cluster and the
rest of the world is exactly two conditions --- and note what is
\emph{absent} from that list. There is no condition of the form
``$T_y$ must not collide with $T_f$'': by Lemma~\ref{lem:ownership} it is
vacuous, because a frozen link's covered ranks live in a non-cluster
interval. Nor is there any condition on $S_y$ versus $S_f$, by
Lemma~\ref{lem:nonconflicts}\,(b). And there is no global condition at
all --- no ordering, no acyclicity --- by Proposition~\ref{prop:noacyclic}.
That is what reduces an exact re-solve to a finite enumeration whose
only external input is the pair of static bitmasks
$\mathrm{Drop}(F)$, $\mathrm{Src}(F)$.

\begin{corollary}[Never worse]
\label{cor:neverworse}
$Y = X_P$ satisfies (C1)--(C4). Hence if $Y^\star$ maximises $\cov(Y)$
subject to (C1)--(C4), then $\cov(Y^\star) \ge \cov(X_P)$ and
\[
  |C|\bigl(F \cup Y^\star\bigr) \;\le\; |C|(X).
\]
An accepted cluster solve can never increase $|C|$, whatever $P$ is.
\end{corollary}

\begin{proof}
$X$ is feasible, so by the converse direction of
Theorem~\ref{thm:compose} (with $Z = X$) the incumbent $X_P$ satisfies
(C1)--(C4); it is therefore in the feasible region of the local
maximisation, which gives $\cov(Y^\star) \ge \cov(X_P)$. Apply the
objective identity.
\end{proof}

\begin{remark}[How the code tests (C3) and (C4)]
\label{rem:avail}
The implementation never materialises $\mathrm{Drop}(F)$ or
$\mathrm{Src}(F)$. It observes that the global arrays
\texttt{removed\_} and \texttt{pin\_count\_} describe $X$, and that
subtracting the cluster's own contribution turns them into a description
of $F$:
\[
\begin{aligned}
  \texttt{kept\_after\_revoke}(r) &\equiv
    \neg\texttt{removed\_}[r] \ \vee\ r \in \textstyle\bigcup_{k \in X_P} T_k
    &&\Longleftrightarrow\quad r \notin \mathrm{Drop}(F), \\
  \texttt{external\_pins}(r) &\equiv
    \texttt{pin\_count\_}[r] - \bigl|\{k \in X_P : r \in S_k\}\bigr|
    &&\Longleftrightarrow\quad
    \bigl|\{f \in F : r \in S_f\}\bigr| .
\end{aligned}
\]
Then \texttt{avail}$(y)$ requires \texttt{kept\_after\_revoke} on every
rank of $S_y$ --- which is (C3) --- and, on every rank of $T_y$, both
\texttt{kept\_after\_revoke} and \texttt{external\_pins} $= 0$; the
latter is (C4), and the former is redundant by
Lemma~\ref{lem:ownership}. Computing this \emph{without} revoking
anything is what allows the solve to be non-mutating: the global state is
touched only when an improvement is actually found
(\S\ref{sec:algo:solve}).
\end{remark}

\section{The algorithm}
\label{sec:algo}

\texttt{tree-dp3} is \texttt{tree-dp} --- preference forest, KEEP/COMPRESS
DP, deepest-first materialization, leftover greedy, Phase~II
unpin/retarget --- followed by the pass described here
(\texttt{run\_phase2\_local\_} in \texttt{compress.hpp}). We describe its
five components in turn.

\subsection{Preprocessing: caches, the dependency graph, and static caps}

\paragraph{Candidate cache.}
For every $c \in \Names$ the static list
$\hat\Links_c = $ \texttt{enumerate\_candidates\_}$(c,\ \text{avail} =
\textsf{false})$ is computed once and sorted by coverage descending, then
$a$ descending. When the pass is invoked from \texttt{tree-dp3} this
cache is handed over from the preference build, so no re-enumeration is
paid. Availability is deliberately \emph{not} used as a filter here: the
cache must be valid across all the mutations the pass performs.

\paragraph{Dependency graph.}
$D$ is built as in Definition~\ref{def:dep}: for each $c$ and each
$k \in \hat\Links_c$, walk $S_k$ and add $\{c, \name(r)\}$. Adjacency
lists are truncated to the degree cap $d_{\max} = 16$
(\texttt{GCSA\_LNS\_DEGREE}). This is a heuristic restriction of the
neighbourhood, not of correctness: by Corollary~\ref{cor:neverworse}, any
$P$ whatsoever yields a sound solve.

\paragraph{Static caps.}
For every name, $\scap(c) := \max_{k \in \hat\Links_c} \cov(k)$ (the head
of the sorted cache, or $0$). Since a feasible $Y$ takes at most one link
per name (C1), $\sum_{c \in P} \scap(c)$ upper-bounds $\cov(Y)$ for every
completion $Y$ of every frozen part. This bound is \emph{static}: it does
not depend on $X$, so it can be compared against the incumbent before any
work is done on the cluster. This is the single most important
performance device in the pass (\S\ref{sec:complexity}).

\subsection{Cluster construction}

Seeds are the names of $\Names$ in decreasing interval size, matching the
greedy Phase~II's bias towards large wins. For a seed $c$, the cluster
$P$ is the BFS ball around $c$ in $D$, truncated at
$q_{\max} = 8$ names (\texttt{GCSA\_LNS\_CLUSTER}); members must have an
interval and a cache entry. Each name is used as a seed at most once per
generation, but may belong to arbitrarily many clusters.

\subsection{The exact solve}
\label{sec:algo:solve}

Algorithm~\ref{alg:solve} states the per-cluster solve. Note the
following design points.

\begin{itemize}
\item \textbf{Incumbent seeding.} \texttt{best} is initialised to the
  incumbent $X_P$ and \texttt{best\_total} to $\cov(X_P)$, so the search
  can only ever return something at least as good
  (Corollary~\ref{cor:neverworse}). Moreover the incumbent link of a
  member is \emph{explicitly appended} to that member's option list if
  it is not already present. This is not redundant: a link produced by
  the retarget loop can carry a composite add $a > a_{\max}$ and
  therefore be absent from the static cache altogether, and even a
  cached link can be cut off by the per-name option cap. Without this
  step ``never worse'' would fail.
\item \textbf{Options.} Member $i$ gets the first $p_{\max} = 12$
  (\texttt{GCSA\_LNS\_OPTS}) entries of $\hat\Links_{c_i}$ with
  $\cov \ge \kappa$ that pass \texttt{avail} (Remark~\ref{rem:avail}),
  i.e.\ conditions (C3)--(C4). Since the cache is coverage-sorted, this
  keeps the most valuable options.
\item \textbf{Branch-and-bound.} Members are ordered by their best option
  coverage, descending, so the search commits to the high-value decisions
  first and the bound bites early. With
  $\mathrm{cap}(i) = \max_{y \in \mathrm{opts}(i)} \cov(y)$ and the
  suffix sums $\mathrm{suf}[k] = \sum_{j \ge k} \mathrm{cap}(\mathrm{ord}[j])$,
  the node at depth $k$ with accumulated coverage $\mathit{tot}$ is
  pruned whenever $\mathit{tot} + \mathrm{suf}[k] \le \mathit{best}$;
  this is valid because at most one link per remaining member may be
  taken, each worth at most its cap.
\item \textbf{The ``no link'' branch.} Every member also has the option of
  taking nothing. It is tried \emph{last}, after all concrete options,
  because it never improves the running total and the bound is most
  effective when a good incumbent has already been established. Its
  presence is what makes the search a search over the full downward-closed
  family $\Feas$ (Lemma~\ref{lem:pairwise}) rather than over perfect
  matchings.
\item \textbf{Internal conflicts.} The pairwise (F2) test against the
  currently live picks uses interval bounds as an $O(1)$ reject
  ($T_{y'} \subseteq [\ell_{c_{y'}}, r_{c_{y'}})$, so the covered ranks are
  scanned only when the source range actually overlaps that interval).
\item \textbf{Non-mutating.} Nothing in the global state is touched until
  a strict improvement is found. On large inputs the overwhelming
  majority of solves end at ``incumbent stands'', and
  \texttt{revoke\_}/\texttt{accept\_} round trips --- which rewrite
  \texttt{removed\_}, \texttt{pin\_count\_} and \texttt{pin\_owner\_} ---
  are the dominant cost if paid speculatively.
\end{itemize}

\begin{algorithm}[t]
\caption{\textsc{SolveCluster}$(P)$ --- exact re-optimization of one cluster}
\label{alg:solve}
\begin{algorithmic}[1]
\Require cluster $P = (c_1,\dots,c_q)$, current assignment $X$, node budget $B$
\Ensure coverage gain $\ge 0$, or $-1$ if the budget was exhausted (state untouched)
\State $\mathit{cur}_i \gets$ link of $c_i$ in $X$ (or none);\quad
       $\mathit{tot}_0 \gets \sum_i \cov(\mathit{cur}_i)$
\State build \textsc{Avail} from $X$ minus the cluster's own links \Comment{(C3)--(C4); Remark~\ref{rem:avail}}
\For{$i \gets 1$ \textbf{to} $q$}
  \State $\mathrm{opts}(i) \gets$ first $p_{\max}$ links of $\hat\Links_{c_i}$
         with $\cov \ge \kappa$ and \textsc{Avail}
  \If{$\mathit{cur}_i$ exists and $\mathit{cur}_i \notin \mathrm{opts}(i)$}
    \State append $\mathit{cur}_i$ to $\mathrm{opts}(i)$ \Comment{keeps the incumbent expressible}
  \EndIf
  \State $\mathrm{cap}(i) \gets \max_{y \in \mathrm{opts}(i)} \cov(y)$ \quad(0 if empty)
\EndFor
\State $\mathrm{ord} \gets$ indices sorted by $\mathrm{cap}(\cdot)$ descending;\quad
       $\mathrm{suf}[k] \gets \sum_{j \ge k} \mathrm{cap}(\mathrm{ord}[j])$
\State $\mathit{best} \gets \mathit{tot}_0$;\quad $Y^\star \gets (\mathit{cur}_i)_i$;\quad
       $\mathit{nodes} \gets 0$
\Procedure{Dfs}{$k$, $\mathit{tot}$, $L$} \Comment{$L$ = links picked so far}
  \State $\mathit{nodes} \gets \mathit{nodes}+1$;
         \textbf{if} $\mathit{nodes} > B$ \textbf{then abort}
  \State \textbf{if} $\mathit{tot} + \mathrm{suf}[k] \le \mathit{best}$ \textbf{then return}
         \Comment{suffix-sum bound}
  \If{$k = q$}
    \State $\mathit{best} \gets \mathit{tot}$;\quad $Y^\star \gets L$;\quad \Return
  \EndIf
  \State $i \gets \mathrm{ord}[k]$
  \ForAll{$y \in \mathrm{opts}(i)$}
    \If{$y$ conflicts with no element of $L$} \Comment{pairwise (F2)}
      \State \Call{Dfs}{$k+1$, $\mathit{tot}+\cov(y)$, $L \cup \{y\}$}
    \EndIf
  \EndFor
  \State \Call{Dfs}{$k+1$, $\mathit{tot}$, $L$} \Comment{leave $c_i$ uncompressed}
\EndProcedure
\State \Call{Dfs}{$0$, $0$, $\emptyset$}
\State \textbf{if} aborted \textbf{then return} $-1$
\State \textbf{if} $\mathit{best} = \mathit{tot}_0$ \textbf{then return} $0$
       \Comment{incumbent stands; nothing mutated}
\State \textbf{for} $i \gets 1$ \textbf{to} $q$: \textsc{Revoke}$(c_i)$;\quad
       \textbf{for} $y \in Y^\star$: \textsc{Accept}$(y)$
\State \Return $\mathit{best} - \mathit{tot}_0$
\end{algorithmic}
\end{algorithm}

\subsection{The driver}

Algorithm~\ref{alg:driver} wraps the solve in the precheck, the fallback,
and the dirty-set generation loop.

\paragraph{The $\sum \scap \le$ incumbent precheck.}
Before a cluster is touched, the pass compares
$\sum_{c \in P} \scap(c)$ against $\cov(X_P)$. If the static upper bound
does not strictly exceed the incumbent, no completion can improve, and
the cluster is skipped without building \textsc{Avail}, without filtering
options, and without entering the search. This is exact --- it never skips
an improvable cluster --- and, empirically, it removes the great majority
of the work (\S\ref{sec:complexity}).

\paragraph{The size-$4$ retry.}
A cluster whose search exhausts the node budget $B = 20000$
(\texttt{GCSA\_LNS\_NODES}) is not abandoned. It is truncated to its
first $q_{\mathrm{fb}} = 4$ BFS members (\texttt{kLnsFallbackCluster}) and
re-solved. Skipping outright would forfeit the improvements a
sub-cluster would still have found, and by Corollary~\ref{cor:neverworse}
the smaller solve is just as sound. Only if the retry also aborts is the
cluster left alone.

\paragraph{Generations.}
When a cluster improves, its members and their $D$-neighbours are marked
dirty for the next generation, since their option sets and their frozen
context have changed. Generations run until the dirty set is empty (a
fixed point) or the shared Phase~II budget (default $100$,
\texttt{kPhase2DefaultIters}) is spent. In practice one to three
generations suffice.

\begin{algorithm}[t]
\caption{\textsc{ClusterLNS} --- the \texttt{tree-dp3} Phase~II pass}
\label{alg:driver}
\begin{algorithmic}[1]
\Require intervals, current assignment $X$ (from the forest DP + retarget loop)
\State build $\hat\Links_c$ for all $c \in \Names$ (reused from the preference build if available)
\State build the dependency graph $D$ (Def.~\ref{def:dep}), truncating degrees to $d_{\max}$
\State $\scap(c) \gets \max_{k \in \hat\Links_c}\cov(k)$ for all $c$
\State $\mathit{dirty} \gets \Names$ sorted by $\isize$ descending
\For{$\mathit{gen} \gets 1$ \textbf{to} $G$ \textbf{while} $\mathit{dirty} \ne \emptyset$}
  \State $\mathit{next} \gets \emptyset$
  \ForAll{seeds $c \in \mathit{dirty}$, each seed used once}
    \State $P \gets$ BFS ball around $c$ in $D$, capped at $q_{\max}$ names
    \If{$\sum_{d \in P} \scap(d) \le \cov(X_P)$}
      \State \textbf{continue} \Comment{static precheck: no completion can improve}
    \EndIf
    \State $g \gets$ \Call{SolveCluster}{$P$}
    \If{$g < 0$ \textbf{and} $|P| > q_{\mathrm{fb}}$}
      \State $P \gets$ first $q_{\mathrm{fb}}$ members of $P$;\quad
             $g \gets$ \Call{SolveCluster}{$P$}
    \EndIf
    \If{$g > 0$}
      \State $\mathit{next} \gets \mathit{next} \cup P \cup N_D(P)$
    \EndIf
  \EndFor
  \State $\mathit{dirty} \gets \mathit{next}$
\EndFor
\end{algorithmic}
\end{algorithm}

\subsection{Correctness}

\begin{theorem}[Soundness and monotonicity]
\label{thm:sound}
Every intermediate assignment produced by Algorithm~\ref{alg:driver} is
feasible, and $|C|$ is non-increasing over the run. Moreover, each
applied solve returns an assignment that is \emph{optimal} among all
completions of its frozen part that are expressible in the member option
lists.
\end{theorem}

\begin{proof}
The pass starts from a feasible $X$ (the output of the forest DP,
leftover greedy and retarget loop, each of which maintains feasibility).
A solve applies $F \cup Y^\star$ where $Y^\star$ satisfies (C1) by
construction (one option per member, options are valid links of member
names), (C2) by the explicit pairwise conflict test, and (C3)--(C4) by
the \textsc{Avail} filter (Remark~\ref{rem:avail}). By
Theorem~\ref{thm:compose}, $F \cup Y^\star$ is feasible; by
Corollary~\ref{cor:neverworse} and the incumbent seeding of line~11 of
Algorithm~\ref{alg:solve}, $\cov(Y^\star) \ge \cov(X_P)$, so $|C|$ does
not increase. Optimality within the option lists is the exhaustiveness of
the DFS: it enumerates every element of
$\prod_i (\mathrm{opts}(i) \cup \{\bot\})$ except those cut by the bound
$\mathit{tot} + \mathrm{suf}[k] \le \mathit{best}$, which discards only
assignments whose total cannot strictly exceed the incumbent best, and
those cut by a pairwise conflict, which are infeasible. If the node
budget is exhausted the solve returns $-1$ and nothing is mutated.
\end{proof}

Note what Theorem~\ref{thm:sound} does \emph{not} claim: global
optimality. The pass is a large-neighbourhood search --- exact within a
neighbourhood, heuristic in the choice of neighbourhoods --- and its
fixed point is a local optimum with respect to the family of BFS balls of
size $\le q_{\max}$. Section~\ref{sec:experiments} measures how far that
falls short of the ILP optimum in practice (usually: not at all, on
instances small enough to check).

\section{Why it composes with, rather than replaces, the retarget loop}
\label{sec:compose-phase2}

Because a cluster solve can never increase $|C|$, it is safe to run
\emph{after} anything. The interesting question is whether it should
\emph{replace} the pre-existing Phase~II unpin/retarget loop, which is
the more expensive of the two. It should not.

The retarget loop does something the cluster solve structurally cannot.
When it accepts a new link $w$ for some name and a dependent $D$ was
sourcing from ranks that $w$ now covers, it \emph{rewrites} $D$ to point
at the corresponding sub-run of $S_w$ with
\[
  a'_D \;=\; a_D + a_w
\]
(\texttt{try\_accept\_with\_retarget\_}: \texttt{d.add += cand.add}).
The synthesized link is a legitimate element of $\Links$ --- (V1) holds
because the two shifts compose --- but it is \emph{not} an element of any
$\hat\Links_c$: \texttt{enumerate\_candidates\_} only ever emits
$a \le a_{\max}$, whereas $a'_D$ may exceed $a_{\max}$, and even when it
does not, the composite source run need not be a maximal
$\mathrm{lcp} \ge a+1$ run. The cluster solve chooses from the static
cache; it can preserve such a link (that is exactly what the incumbent
appending of Algorithm~\ref{alg:solve}, line~6, is for) but it can never
invent one.

On long repeats this matters a great deal, because composite adds are how
a deep chain of repeated motifs gets collapsed onto a single stored
prefix. Running the cluster solve \emph{instead of} the retarget loop
(\texttt{GCSA\_LNS\_ONLY=1}) loses ground on exactly those inputs: over
the $139$-configuration suite at $a_{\max} = 8$,

\begin{center}
\begin{tabular}{@{}lrr@{}}
\toprule
Phase~II configuration & total $|C|$ & vs.\ \texttt{tree-dp} \\
\midrule
retarget loop only (\texttt{tree-dp})           & $252397$ & --- \\
cluster solve only (\texttt{GCSA\_LNS\_ONLY=1}) & $253675$ & $+1278$ \\
retarget loop then cluster solve (\texttt{tree-dp3}) & $252394$ & $-3$ \\
\bottomrule
\end{tabular}
\end{center}

The two passes are therefore complementary, not competing: the retarget
loop enlarges the link universe along composite-add chains, and the
cluster solve makes exact choices within the static universe. The
composition is safe in one direction only --- the solve after the loop,
never in place of it.

\section{Complexity}
\label{sec:complexity}

\paragraph{Preprocessing.}
Building $D$ costs one pass over the source runs of all cached
candidates, $O\bigl(\sum_{c}\sum_{k \in \hat\Links_c} |S_k|\bigr)$, plus
the linear-scan deduplication of adjacency lists, which is $O(d_{\max})$
per insertion after the degree cap. Static caps are $O(|\Names|)$, since
the caches are already coverage-sorted.

\paragraph{One cluster.}
A cluster has $q \le q_{\max}$ members with at most $p_{\max}$ options
each, plus the ``no link'' branch, so the raw search tree has at most
\[
  (p_{\max} + 1)^{q_{\max}} \;=\; 13^8 \;\approx\; 8.2 \times 10^8
\]
leaves in the worst case --- which is why the node budget $B$ exists at
all. Each visited node performs $O(q)$ conflict tests, each $O(1)$
except on a genuine source/interval overlap, where it is $O(\cov)$. In
practice the suffix-sum bound and the descending-capacity member order
collapse this by orders of magnitude: the very first root-to-leaf path
takes the highest-capacity option of every member, which typically
establishes a bound close to $\mathrm{suf}[0]$ and prunes almost
everything else. The budget therefore caps the pathological cases rather
than the typical ones; the fallback then re-solves at $q = 4$, i.e.\ at
most $13^4 \approx 2.9 \times 10^4$ leaves.

\paragraph{The pass.}
Each generation seeds at most $|\Names|$ clusters; cluster construction is
$O(q_{\max} d_{\max})$ and the precheck is $O(q_{\max})$. So the
per-generation cost is
\[
  O\bigl(|\Names| \cdot q_{\max} d_{\max}\bigr)
  \;+\;
  \sum_{\text{clusters not skipped}} O(B \cdot q_{\max}),
\]
and the number of generations is bounded by the Phase~II budget $G$.
Because a generation's dirty set is confined to the neighbourhoods of the
clusters that actually improved, later generations are far cheaper than
the first.

\paragraph{The precheck in practice.}
The second term above is what the $\sum \scap \le$ incumbent test
attacks, and it attacks it very effectively. On
\texttt{r300-3000.fasta} ($n = 9\,000\,000$) with shape
\texttt{\#\#\#\#.\#\#\#\#} and $a_{\max} = 8$ there are $24034$ names with
$|I_c| > 1$, hence $24034$ seeded clusters; the precheck skips
\[
  22168 \ \text{of}\ 24034 \quad (92.2\,\%),
\]
leaving $1866$ actual solves, of which none improved (with $a_{\max}=8$
the retarget loop has already saturated this instance) and none aborted.
The whole pass costs $0.61$\,s ($0.20$\,s of which is building $D$) on top
of a $23.4$\,s retarget loop:

\begin{center}
\begin{tabular}{@{}lrrr@{}}
\toprule
algorithm & Phase~II & total compress & $|C|$ \\
\midrule
\texttt{tree-dp}  & $23.6$\,s & $24.0$\,s & $1919406$ \\
\texttt{tree-dp3} & $24.0$\,s & $24.5$\,s & $1919406$ \\
\bottomrule
\end{tabular}
\end{center}

so the exact pass is at parity with \texttt{tree-dp} on a nine-megabase
input --- it adds about $2\,\%$ to the compression time. Without the
precheck, and with a mutating solve that pays
\texttt{revoke\_}/\texttt{accept\_} on every cluster rather than only on
improvement, the same pass was reported to take about $352$\,s; we were
not able to reproduce that figure with the current code (the two
optimizations are not switchable), and record it here as the historical
motivation for both devices rather than as a measurement.

\section{Experimental results}
\label{sec:experiments}

\subsection{The $139$-configuration suite}

\texttt{compare\_algos} runs five hand instances, four short repetitive
instances, and a weight-$30$ suite of $13$ shapes $\times$ $10$
repetition counts of a random $200$\,bp motif ($139$ configurations in
total), verifying every result against brute force and against a
serialized round-trip. Totals of $|C|$ over all $139$ configurations:

\begin{center}
\begin{tabular}{@{}lrrrrr@{}}
\toprule
$a_{\max}$ & greedy & dep-order & \texttt{tree-dp} & \texttt{tree-dp2} & \texttt{tree-dp3} \\
\midrule
$8$  & $259518$ & $252357$ & $252397$ & $256197$ & $\mathbf{252394}$ \\
$16$ & $255785$ & $252175$ & $252329$ & $255370$ & $\mathbf{213078}$ \\
$32$ & $253709$ & $252155$ & $252177$ & $252917$ & $\mathbf{225814}$ \\
\bottomrule
\end{tabular}
\end{center}

Per-configuration, \texttt{tree-dp3} is never worse than \texttt{tree-dp}
at any $a_{\max}$ --- as Corollary~\ref{cor:neverworse} requires, since
the two share everything up to the final pass:

\begin{center}
\begin{tabular}{@{}lrrrr@{}}
\toprule
$a_{\max}$ & better than \texttt{tree-dp} & equal & worse & reduction in total $|C|$ \\
\midrule
$8$  & $2$   & $137$ & $0$ & $0.001\,\%$ \\
$16$ & $122$ & $17$  & $0$ & $15.56\,\%$ \\
$32$ & $122$ & $17$  & $0$ & $10.46\,\%$ \\
\bottomrule
\end{tabular}
\end{center}

\subsection{The counterexamples and the ILP optimum}

On small instances the exact optimum is available from
\texttt{ilp\_baseline}, which enumerates the full link universe of
Definition~\ref{def:link} (all sub-runs, no maximality or lcp filter) and
solves the set-packing ILP with \texttt{cbc}. The two counterexamples of
\S\ref{sec:ce1}--\ref{sec:ce2}, and the hardest case of the fixed
suite:

\begin{center}
\begin{tabular}{@{}llrrrrrr@{}}
\toprule
instance & shape & opt & greedy & dep-order & \texttt{tree-dp} & \texttt{tree-dp2} & \texttt{tree-dp3} \\
\midrule
$(\mathtt{AC})^6$              & \texttt{\#.\#}  & $9$  & $10$ & $9$  & $10$ & $10$ & $\mathbf{9}$ \\
$(\mathtt{ACGT})^5$            & \texttt{\#.\#}  & $11$ & $14$ & $11$ & $13$ & $13$ & $\mathbf{11}$ \\
$(\mathtt{GCCTTTAAAG})^4$      & \texttt{\#..\#} & $22$ & $25$ & $27$ & $27$ & $27$ & $\mathbf{22}$ \\
note text ($n=40$)             & \texttt{\#\#}   & $19$ & $20$ & $20$ & $20$ & $20$ & $20$ \\
\bottomrule
\end{tabular}
\end{center}

Over the whole ten-case fixed suite (\texttt{test\_ilp\_cases.sh}),
\texttt{tree-dp3} attains the ILP optimum on $9$ of $10$ instances; the
sole exception is the last row above, where every heuristic stores one
position too many. The $(\mathtt{GCCTTTAAAG})^4$ / \texttt{\#..\#} case is
the most striking: the optimum needs six simultaneous links, and the
cluster solve finds all six, cutting $|C|$ from \texttt{tree-dp}'s $27$ to
$22$ --- an $18.5\,\%$ reduction on that instance, and also better than
greedy's $25$, which no other heuristic achieves.

\subsection{On the ILP being a lower bound}
\label{sec:ilpcaveat}

\texttt{ilp\_baseline} reports \texttt{bounds all algos}, i.e.\ whether
the computed optimum is $\le$ every heuristic's $|C|$. This is checked,
not proven. The universe it enumerates is
$\{k \in \Links : a_k \le a_{\max}\}$, and by
\S\ref{sec:compose-phase2} the retarget loop can synthesize links with
$a > a_{\max}$; such a link is decodable and would be feasible in a
larger universe, so in principle a heuristic could reach a solution the
ILP never considered. We are not aware of an instance where this
happens, and the randomized stress harness (\texttt{test\_ilp\_stress.sh},
$410$ instances of length $10$--$39$ over two alphabets and six shapes)
reports no violation --- but the claim is empirical. Closing it properly
would require enumerating links with $a$ up to the largest composite add
any heuristic can build, which is bounded only by the text length.

\section{Discussion}
\label{sec:discussion}

\subsection{At the default $a_{\max}$, the exact solve is starved}

The most important number in \S\ref{sec:experiments} is the smallest one.
At $a_{\max} = 8$, exact re-optimization over bounded dependency
neighbourhoods --- a genuinely exact procedure, applied to every name in
the instance --- buys $3$ stored positions across $139$ configurations,
all of them on the two hand-built counterexamples ($10 \to 9$ and
$13 \to 11$). On all $137$ real configurations the answer is unchanged.

This is not a failure of the neighbourhood size. At $a_{\max} = 8$ on the
large repetitive instances, the precheck reports that for over $90\,\%$
of clusters the \emph{static} bound $\sum_{c \in P} \scap(c)$ does not
even exceed what the incumbent already achieves: there is no better
assignment to find, no matter how large a neighbourhood one is willing to
solve, because the candidate lists do not contain anything better. What
binds is the \emph{candidate universe}, not the search.

Raising $a_{\max}$ to $16$ enlarges that universe and the same algorithm
becomes worth $15.5\,\%$ of total $|C|$, improving $122$ of $139$
configurations. It is worth being precise about the division of labour
here: greedy, dep-order, \texttt{tree-dp} and \texttt{tree-dp2} \emph{all}
barely move between $a_{\max} = 8$ and $a_{\max} = 16$ (totals change by
under $1.5\,\%$), because none of them can exploit the extra candidates
--- they commit to one preferred link per name and the preferred link
rarely changes. The exact solve can, because it evaluates the whole
option list against its neighbours. The lesson generalises: enriching the
candidate universe pays only if the selection procedure is strong enough
to use it, and strengthening the selection procedure pays only if the
universe is rich enough to reward it. \texttt{tree-dp3} at
$a_{\max} = 16$ is the first configuration in which both halves are in
place.

\subsection{The $a_{\max} = 32$ regression is a neighbourhood-size artifact}

Going from $a_{\max} = 16$ to $32$ makes \texttt{tree-dp3} worse
($213078 \to 225814$) while leaving every other algorithm essentially
unchanged. Since larger $a_{\max}$ strictly enlarges the candidate
universe, this must be an artifact of the caps rather than of the
problem. We measured which cap:

\begin{center}
\small
\begin{tabular}{@{}lrrrrr@{}}
\toprule
suite configuration ($\times 100$ reps, $a_{\max}=32$)
  & $q_{\max}=8$ & aborts & $q_{\max}=12$ & aborts & ($a_{\max}=16$, $q_{\max}=8$) \\
\midrule
\texttt{5+g4+20+g4+5}  & $3270$ & $0$ & $2796$ & $25$ & $2717$ \\
\texttt{10+g8+20}      & $3270$ & $0$ & $2890$ & $28$ & $2807$ \\
\texttt{5+g2+20+g2+5}  & $3063$ & $1$ & $2679$ & $29$ & $2881$ \\
\bottomrule
\end{tabular}
\end{center}

The node budget is not the binding constraint: at $q_{\max} = 8$ the
abort count is $0$ or $1$, and raising \texttt{GCSA\_LNS\_NODES} by a
factor of $20$ changes nothing at all. Raising the \emph{cluster} cap
from $8$ to $12$ recovers most of the loss --- and on one of the three
configurations overshoots the $a_{\max}=16$ result. The mechanism is
straightforward once stated: a larger $a_{\max}$ gives each name
candidates sourcing from more distinct intervals, so the dependency
graph densifies, and a BFS ball of a fixed $8$ names covers a
progressively smaller fraction of each name's true dependency
neighbourhood. More of the relevant structure ends up frozen at the
cluster boundary, and conditions (C3)--(C4) of
Theorem~\ref{thm:compose} bind harder. The regression is thus a tuning
artifact of $q_{\max}$, and the natural fix is to scale $q_{\max}$ with
the observed dependency degree rather than fixing it at $8$.

(An earlier account attributed this regression to node-cap aborts. The
measurements above do not support that: at the default cluster cap the
aborts are essentially absent, and the node budget is not the active
constraint.)

\subsection{What remains}

Three directions follow directly from the analysis above.

\begin{enumerate}
\item \textbf{Cluster sizing.} $q_{\max}$ is a fixed constant chosen for
  the $a_{\max}=8$ regime. Given a per-instance measurement of the
  dependency degree it could be chosen adaptively, with the node budget
  --- which is currently slack --- absorbing the variance.
\item \textbf{Bringing composite adds into the universe.} The retarget
  loop's composite links are exactly the ones the static cache lacks
  (\S\ref{sec:compose-phase2}). Emitting them into $\hat\Links_c$, rather
  than only preserving them when they happen to be incumbent, would let
  the exact solve reason about them --- and would also close the ILP
  lower-bound caveat of \S\ref{sec:ilpcaveat}.
\item \textbf{Beyond BFS balls.} The neighbourhood family is currently
  ``bounded balls around each name''. Since the algorithm is exact within
  whatever family it is given, a family chosen from the structure of the
  incumbent --- for instance the tight cycles of the preference graph,
  which are precisely what \S\ref{sec:ce1} and \S\ref{sec:ce2} show
  \texttt{tree-dp} mishandles --- might buy more per unit of work than
  uniform balls.
\end{enumerate}

What the present note establishes is the structural point underneath all
three: because the compressed array stores literal positions, the
feasible region of differential compression is a plain set-packing
region with pairwise constraints and no acyclicity condition, and
because a link's covered ranks lie in its own interval, any set of names
can be carved out and re-optimized exactly against a two-condition
interface. Exactness on a neighbourhood is therefore cheap and always
safe. Whether it is \emph{useful} is entirely a question of how rich the
candidate universe is, and at the default $a_{\max}$ it is not yet rich
enough.

\end{document}
